\chapter{Thunderclap Headache}

\section*{Definition}
A thunderclap headache is a sudden, severe headache that reaches maximal intensity within seconds to one minute. It may be the first sign of bleeding around the brain, cerebral venous thrombosis, arterial dissection, reversible cerebral vasoconstriction syndrome, or another vascular emergency. A new thunderclap headache requires immediate medical assessment even if the pain improves.

\section*{When to See a Doctor}
Seek emergency care immediately for any first or unusual thunderclap headache. Emergency symptoms include neck stiffness, vomiting, fainting, confusion, seizure, fever, weakness, numbness, speech difficulty, visual loss, unequal pupils, pregnancy or recent childbirth, head trauma, or headache triggered by exertion or sexual activity. Do not drive yourself.

\section*{How It Develops}
The abrupt pain reflects sudden irritation or stretching of pain-sensitive structures in the head, neck, or meninges. Subarachnoid hemorrhage from a ruptured aneurysm is a critical cause. Other causes include reversible cerebral vasoconstriction, cervical artery dissection, cerebral venous thrombosis, intracranial hemorrhage, pituitary apoplexy, hypertensive emergency, meningitis, and spontaneous intracranial hypotension. Primary thunderclap headache is diagnosed only after dangerous secondary causes have been excluded.

\section*{Causes}
\begin{itemize}
  \item Subarachnoid, intracerebral, or subdural hemorrhage
  \item Reversible cerebral vasoconstriction syndrome, sometimes associated with pregnancy, the postpartum period, vasoactive drugs, or intense exertion
  \item Cervical artery dissection, cerebral venous sinus thrombosis, or ischemic stroke
  \item Meningitis, encephalitis, hypertensive emergency, or pituitary apoplexy
  \item Head or neck trauma, sexual activity, heavy exertion, coughing, or straining
  \item Primary thunderclap headache after appropriate investigation excludes secondary disease
\end{itemize}

\section*{Related Symptoms}
\begin{itemize}
  \item Neck pain or stiffness, nausea, vomiting, photophobia, or loss of consciousness
  \item Weakness, numbness, facial droop, speech or language difficulty, imbalance, or seizure
  \item Fever, rash, visual disturbance, or severe pain extending into the neck or back
\end{itemize}
\textbf{Important:} A normal appearance between episodes does not safely exclude subarachnoid hemorrhage or reversible cerebral vasoconstriction.

\section*{Medical Evaluation}
\begin{itemize}
  \item Clinicians establish the time to peak intensity, associated neurologic symptoms, triggers, prior headache pattern, and vascular or pregnancy risk.
  \item Examination includes vital signs, funduscopic and neurologic examination, meningismus, and assessment for trauma or arterial dissection.
  \item Noncontrast head CT is performed urgently; CT angiography may assess aneurysm, dissection, or vasoconstriction.
  \item If initial imaging is nondiagnostic and suspicion remains, lumbar puncture or MRI/MRA/MRV may be required, guided by timing and specialist judgment.
\end{itemize}

\section*{Treatment Options}
Treatment depends on the cause and may include neurosurgical or endovascular management of hemorrhage, blood-pressure control, treatment of infection, anticoagulation for selected venous thrombosis, or specialist management of vasoconstriction. Analgesia and antiemetics are supportive measures and must not delay diagnostic evaluation.

\section*{Self-Care and Home Management}
Do not wait for home treatment to work and do not repeatedly take analgesics to mask the symptom. Note the exact onset time, triggers, associated symptoms, recent infections, pregnancy or postpartum status, medications, recreational drugs, and trauma. Call emergency services and remain with someone until assessed.

\section*{References}
\begin{thebibliography}{99}
\bibitem{thunderclap_headache:ref1} Perry JJ, et al. Sensitivity of computed tomography performed within six hours of headache onset for diagnosis of subarachnoid haemorrhage. \textit{BMJ}. 2011;343:d4277.
\bibitem{thunderclap_headache:ref2} Connolly ES Jr, et al. Guidelines for the management of aneurysmal subarachnoid hemorrhage. \textit{Stroke}. 2023;54:e314--e370.
\bibitem{thunderclap_headache:ref3} Miller TR, Shivashankar R, Chung LS, Rispoli JV. Reversible cerebral vasoconstriction syndrome, part 1. \textit{AJNR Am J Neuroradiol}. 2015;36(8):1392--1399.
\end{thebibliography}
